We evaluate four research questions (RQs). RQ1 compares classification performance, RQ2 examines the construction mechanisms, RQ3 tests memory reuse across student architectures and relative to training on all observed data, and RQ4 separates construction costs from downstream training costs.

\subsection{Experimental Setup}

\noindent\textbf{Datasets and continual protocol.}
\begin{table}[t]
\centering
\footnotesize
\setlength{\tabcolsep}{3pt}
\caption{Construction pools, final memory sizes, and class arrival schedules used in the main evaluation.}
\label{tab:datasets}
\begin{tabular*}{\columnwidth}{@{\extracolsep{\fill}}ccccc@{}}
\toprule
\multicolumn{1}{c}{Dataset} &
\multicolumn{1}{c}{\shortstack{Construction nodes}} &
\multicolumn{1}{c}{\shortstack{Memory items}} &
\multicolumn{1}{c}{Stages} &
\multicolumn{1}{c}{\shortstack{Class schedule}} \\
\midrule
ArXiv     & 118,539 & 5,927 & 3 & 13/14/13 \\
SportsFit & 121,137 & 6,057 & 5 & 5/2/2/2/2  \\
BookHis   & 29,086  & 1,454 & 5 & 5/2/2/2/1  \\
Cora      & 1,895   & 95    & 3 & 2/2/3      \\
Citeseer  & 2,229   & 111   & 3 & 2/2/2      \\
\bottomrule
\end{tabular*}
\end{table}

We evaluate \method{} on ArXiv, SportsFit, and BookHis, covering citation
and product networks, and include Cora and Citeseer as smaller graphs. The
standard citation sources are described in~\cite{tag_benchmark,cora,citeseer},
and the academic network subsets are related to the AMiner collection
~\cite{arnetminer}. Product-network construction is consistent with the
textual product relations studied in~\cite{amazon_graph,pin_sage}.
\autoref{tab:datasets} gives the construction pools and class arrival
schedules. The main class-incremental comparison uses three ArXiv stages with
13/14/13 classes. SportsFit and BookHis use five stages with the schedules in
the table. Cora and Citeseer use three stages because they contain only seven
and six classes. All methods retain a fixed \(5\%\) of the observed
construction pool. The same class order determines each method's current
nodes and memory allocation. The ArXiv comparison in the main table uses this
three-stage schedule.

We also evaluate five stages ordered by publication year on ArXiv and two
additional citation graphs, ACM and DBLP. Classes can recur, so this is a
different protocol from the disjoint class-incremental schedules above.
Construction uses the current arrival group and its induced edges. Cumulative
prediction uses all nodes observed up to the current stage and the edges
between them. ArXiv uses year cutoffs of 2015, 2016, 2017, 2018, and 2020.
ACM and DBLP use quantiles of training publication years, keeping papers from
the same year together. The year order is fixed across construction seeds.
These streams capture publication chronology, with edge availability
determined by the arrival of both endpoints. The final cumulative results
and their current-arrival diagnostic are reported separately in
\autoref{sec:main_results} and the compact appendix. Dataset details and
numerical settings are given above. Dynamic graph studies provide the broader
context for this chronological protocol~\cite{evolvegcn,tgn,dynamic_survey}.

\noindent\textbf{Baselines.}
CaT~\cite{cat} uses distribution matching, and GECC~\cite{gecc} uses
clustering of propagated features with hard assignments. Both attach source
text to their numerical memories. Random, K-center, cluster medoid, and
herding select real nodes with their texts. K-center selection follows the
core-set view of active learning~\cite{kcenter}, medoid selection follows
the standard cluster representative formulation~\cite{kmedoids}, and
herding is a common exemplar rule in incremental learning~\cite{icarl}.
GCond-x~\cite{gcond} uses the simplified graph convolution model~\cite{sgc}
and a graph condensation objective, while SFGC~\cite{sfgc} adapts trajectory
matching to graph-free data. TAGSAM~\cite{tagsam} retains text selection and
response matching with a reduced construction budget. STAD~\cite{stad} uses
coarsening, self-training, and a local text generator. The method names in
our tables denote adaptations to common access, memory, and student
requirements. Implementation and budget changes are specified in the table
captions and the surrounding text.

\noindent\textbf{Metrics and training.}
We report the macro-averaged F1 score (Macro-F1), which weights classes
equally, and Accuracy, both as percentages. The F1 score is the harmonic
mean of precision and recall. Main comparisons use the same historical test
partition, with construction and development nodes excluded. This partition
has been inspected in earlier experiments and is not a fresh blind test. The
student uses a two-layer graph convolutional network (GCN)~\cite{gcn}, class
prompt supervision, and a frozen MiniLM text encoder~\cite{minilm}. It trains
for 2,048 updates with batch size 24. Architecture transfer replaces the
graph network with a two-layer GraphSAGE model~\cite{graphsage} while reusing
the same memories. Construction settings, rate sensitivity, and the
historical memory policy follow the definitions above.

A class order controls class arrival, a construction seed controls memory
construction, and a student seed controls downstream training. The target
design crosses three class orders with three construction seeds, giving nine
memories per method and dataset. Each memory is evaluated with three student
seeds. We average its three student scores, then report the mean and sample
standard deviation (SD) across the nine memory scores. SD is used throughout
the paper for this sample standard deviation. All cells in
\autoref{tab:main_comparison} use these nine-memory estimates. Paired 95\%
confidence intervals use differences between matched memory scores. These
intervals summarize the observed crossed design, whose orders and seeds are
shared across runs. Chronological and mechanism experiments use one fixed
order, three construction seeds, and three student seeds per memory unless
stated otherwise.
